\scanpage{5}

La propriété essentielle de $K(C,\delta)$ \deleted{\unsure{a}} est résumée \unsure{dans} :

a) $\tau_i:C_i\unsure{-|\to}C_{i+1}$ est \deleted{\unsure{préexact pour $\delta$|préexact}} \inserted{$\delta_i$-préexact}, i.e.

\[
a')\left\{\begin{aligned}
a'_1)&\ \deleted{\text{Unicité de ttes flèches }z\rightrightarrowsx}\ \text{Pour }(x\xleftarrow{p_\lambda}x_\lambda)\in d'_i,\text{ et ttes doubles}\\
&\text{flèches }z\overset{u,v}{\rightrightarrows}x\text{ telles que }p_\lambda\circ u=p_\lambda\circ v\quad\forall\lambda,\text{ on a }\tau_i(u)=\tau_i(v)\\
a'_2)&\ \text{Pour tte }\deleted{\unsure{\text{sys}}}\ \text{\uline{augmentation} }(z\xrightarrow{q_\lambda}x_\lambda)\text{ du syst. proj., il existe}\\
&\tau_i(z)\xrightarrow{q}\tau_i(x)\text{ telle que }\tau_i(p_\lambda)\circ q=\tau_i(q_\lambda)\quad\forall\lambda
\end{aligned}\right.
\]

a$''$) \deleted{\unsure{comme}} dual de a\unsure{′} ), avec $d''_i$ au lieu de $d'_i$.

\deleted{b) Tout diagramme $\unsure{d'_i|d'_{i+1}},\unsure{\Theta|\supset}\in d'\quad\delta'_{i+1}\supset\tau_i(\delta'_i)$}

\[
b)\left\{\begin{aligned}
b')&\ \forall\text{ diagramme }\inserted{\Theta'}\text{ de type }d'\text{ dans } C_i,\ \exists\text{ une }\deleted{\unsure{\text{augmen.}}}\inserted{\text{augmentation}}\text{ de}\\
&\tau_i(\Theta')\text{ qui est }\in\delta'_{i+1}\\
b'')&\ \forall\text{ diagramme }\Theta''\text{ de type }d''\text{ dans }C_i,\ \exists\text{ une augmentation de }\tau_i(\Theta'')\inserted{\text{ dans }\delta''_{i+1}}.
\end{aligned}\right.
\]

\[
c)\left\{\begin{aligned}
&\delta'_{i+1}\ \deleted{=}\ \inserted{\text{est}}\ \deleted{\unsure{\text{formé}}}\ \tau_i(\delta'_i)\cup\text{ diagr. }\deleted{\unsure{\in\delta'_{i+1}}}\text{ }\unsure{\text{munis}|\text{nouveaux}}\text{ de type }d'\text{ construits dans }b',\\
&\delta''_{i+1}=\tau_i(\delta''_i)\cup\text{ diagrammes de type }d''.\\
&\deleted{\text{tout élément de }\delta'_{i+1}\ (\delta''_{i+1})\text{ de type }\Theta'\ \unsure{\notin}d'\ (\Theta''\notin d'')\text{ est dans }\tau_i(\delta'_i),\ (\tau_i(\delta''_i))}
\end{aligned}\right.
\]

% margin note written sideways
\margin{de c) on conclut par récurrence transfinie\\
$(c^*)\left\{\begin{aligned}\delta'_i&=\alpha_i(\delta')\cup\deleted{\unsure{\text{diagr. de }\delta'_i}}\ \unsure{\text{images dans }C_i\text{ des}|\text{images des}}\ \deleted{\text{diagrammes}}\inserted{\text{cônes proj}}\text{ de type }d'\deleted{\unsure{\text{ contenus}}}\\
\delta''_i&=\alpha_i(\delta'')\cup\deleted{\unsure{\text{diagr. de }\delta''_i}}\ \unsure{\text{images dans }C_i\text{ des cônes ind.}}\text{ de type }d''.
\end{aligned}\right.$}

d) $\forall$ Catégorie $E\in\mathrm{Ob}\,\Delta$, le foncteur $\varphi\longmapsto\varphi\circ\tau_i$ induit une équivalence
\[
\underline{\mathrm{Hom}}_{\delta_{i+1}}(C_{i+1},E)\longrightarrow\underline{\mathrm{Hom}}_{\delta_i}(C_i,E)
\]

\deleted{La} La construction de $K:(C,\delta)- K(C,\delta)$ sera donnée plus bas. On \deleted{\unsure{doit|peut}} \inserted{va} montrer comment elle permet de terminer.

\deleted{De c) on conclut par récurrence transfinie,}

De d), \uline{on déduit par récurrence transfinie}
\[
(d^*)\qquad\boxed{\underline{\mathrm{Hom}}_{\delta_i}(C_i,E)\xrightarrow{\sim}\underline{\mathrm{Hom}}_\delta(C,E)}\qquad \varphi\longmapsto\varphi\circ\alpha_i
\]

Donc on a terminé si on montre qu'il existe $\omega$ \inserted{ordinal} tel que
\[
\left\{\begin{aligned}
1^\circ)&\ C_{\deleted{\unsure{i}}\omega}\in\mathrm{Ob}\,\Delta\qquad 2^\circ)\ \alpha_\omega:C\to C_\omega\ \deleted{\text{est}}\ \delta\text{ exact}\\
3^\circ)&\ \underline{\mathrm{Hom}}_{\delta_\omega}(C_\omega,E)=\underline{\mathrm{Hom}}_{\Delta,\tau_\omega(\delta)}(C_\omega,E)
\end{aligned}\right.
\]

\deleted{\unsure{$2^\circ$) Soit}} Soit pour ceci $\pi$ un cardinal tel que
\[
\pi>\mathrm{card}\,\mathrm{Fl}\,\Theta,\ \deleted{\unsure{\text{card Ob}}}\left\{\begin{aligned}
&\text{si }\Theta\ \deleted{\unsure{\in d'\text{ ou}}}\ \inserted{\in d'\text{ ou }\Theta\in d''}\\
&\text{ou }\Theta\ \deleted{\unsure{\in}}\text{ indexe un }\in\delta'\text{ ou un }\in\delta''
\end{aligned}\right.
\]
et $\omega$ le plus petit ordinal de cardinal $=\pi$. Donc l'ens $S_\omega$ des

% archive sheet number: 4
